% !TEX program = pdflatex
\documentclass[12pt]{article}
\usepackage[utf8]{inputenc}
\usepackage{graphicx}
\usepackage{hyperref}
\usepackage{amsmath}
\usepackage{amsfonts}
\usepackage{booktabs}
\usepackage[numbers]{natbib}
\usepackage[a4paper, margin=1in]{geometry}

\title{Utilizing Topic Modeling for Music Recommendation}
\author{Le Tran Minh Khoa}

\begin{document}
\maketitle

\begin{abstract}
    Recommendtion systems have been rising in many aspects of our daily lives as a means to filter and suggest contents for the customers based on their preferences, with its appearances in so many industries.In music streaming platforms like Spotify and Apple Music, personalized suggestions help users discover new tracks, increasing engagement and retention. However, popular algorithms often rely on collaborative filtering and audio metadata, which can amplify popularity bias, leading to homogeneous recommendations dominated by mainstream artists. In this work, we explore the use of lyrics-based recommendation systems leveraging topic modeling techniques to provide diverse and coherent music suggestions. We compare classical methods like Latent Dirichlet Allocation (LDA) with neural approaches such as BERTopic and Contextualized Topic Models (CTM). Our findings indicate that CTM effectively captures lyrical themes, producing recommendations that reduce popularity bias while maintaining thematic relevance.
\end{abstract}

\section{Introduction}
Recommendation systems are integral to modern digital services, reducing
information overload by filtering content based on user preferences. In the
domain of music, recommender systems have been widely impplemented to enhance
user experience by suggesting tracks that align with individual tastes, which,
otherwise, would be difficult to discover in the vast catalog of available
music. There have been many approaches and efforts in the research for the
music recommendation system, which have been evolving non-stop throughout the
years. In the beginning, the recommendation systems were mostly based on
collaborative filtering, which relies on user-item interactions to suggest
items that similar users liked. The idea is to recommend items that the other
users with similar tastes to the target user to that target user, or proposing
the items similar to the items already favored by the target user. However, one
notable problem with this approach is the cold-start problem, where new users
or items lack sufficient interaction data to generate meaningful
recommendations.Then, we move on to some methods that are based on deep
learning methods Considering this issue, using convolutional neural networks to
extract audio features and metadata such as genre, artist popularity, and
release date has come up as a viable solution to enhance the recommendation.
One of such work proposed integrating the use of convolutional neural networks
to categorize music based on the audio features while the collaborative
filtering algorithm is used to furnish the recommendation based on the
categorized music data \cite{mrcnn}. Then, there comes the approach using
recurrent neural networks to capture the sequential nature of music listening
behavior, allowing for more personalized recommendations based on user
listening patterns, since the listening behavior of the users is not static,
but rather dynamic and changes over time.The CoSeRNN model, for example, uses
RNNs to generate session-aware embeddings that reflect user preferences in
real-time \cite{hansen2020contextual}. Furthermore, with the rise of attention
mechanisms and transformer architectures, with the huge advantage of being ablt
to capture long-range context and consider both the past and the future
interactions of the users - in other words, bilaterally instead of only looking
at the past interactions of the user (one-sided) like the RNNs.

Lyrics-based recommendation offers a complementary signal, focusing on thematic
content rather than popularity. Unlike audio features or metadata, topic
modeling of lyrics captures semantic patterns that reflect an artist's style
and themes. This approach reduces dependence on external metadata and mitigates
the tendency to over-favor popular artists. Our work investigates the efficacy
of topic modeling for music recommendation, comparing classical and neural
methods to deliver diverse, coherent suggestions. Reference
\cite{yadav2020bidirectional} presents a sequential music recommender system
using BERT, called Bert4Vec, which captures sequential patterns in user
listening behavior for improved music recommendations. The model is designed to
understand the context and dependencies in music listening sessions, providing
more accurate and personalized recommendations. Another approach is to use
graph neural networks to model the relationships between the users and the song
items, which can be represented by connections between users and items in a
graph structure, allowing for more nuanced recommendations based on the network
of interactions.

If we look at another perspective in terms of the approach, we could also
categorize the recomendations into content-based, which is the use of the music
various metadata such as melody, rythm,....; context-based, which might
encompass details such as playlists, gender, age, location, etc...;
emtion-based, in which we try to capture the emotions conveyed in the songs by
various methods, such as sentiment analysis or emotion recognition models.
Lately, there also emerge reinforcement-based recommendation systems, which
learn to recommend items by maximizing long-term user engagement through
trial-and-error interactions, rather than putting the recommendation system's
goal in a prediction or classification task.

For further reference about the ongoinng status of the research in this music
recommendation system field, please refer to the survey paper
\cite{musicrecsurvey}.

Amidst all of these approaches, this work chooses to focus on analyzing the
lyrics of the song, which is a rich source of information that can be used to
understand the underlying themes, emotions, and topics of a song. By analyzing
the lyrics, we can gain insights into the artist's style, the song's genre, and
the overall mood of the song, which can be categorized as content-based
recommendation. Although in music, the audio features might be the more
prominent factor for many people, I believe that the all the good songs have
something in common, which is the lyrics, which will help to transfer the
artist's message and the inner feelings to the audiances. I also believe that
while the audio part of the song might be what attracts the listeners at first,
it is the lyrics that will hook the listeners and make them stay, which makes
the song memorable, meaningful, and standing through the test of time.
Therefore, I believe that the lyrics-based recommendation system is a viable
approach to recommend songs to the users, which can be used in conjunction with
other approaches such as audio features or metadata. In this work, I will focus
on the use of topic modeling to first analyze the lyrics of the songs and
extract out the possible potential topics, or the general themes of the songs,
and then recommend songs which are similar to the query song in terms of the
topics. By doing this way, this work hopes to achieve to provide a diverse
recommendation system which can also mitigate the effect of popularity bias,
which is a well-known problem in the music recommendation system: by not
utilizing the audio features and just simply based on the lyrics, the
recommendation might be able to avoid songs that might sound similar, in which
many cases belong to the same artist, reducing the discoverability of new
artists and songs; by not utilizing the metadata, the recommendation might be
able to avoid songs that are popular but not necessarily good, which might be a
result of the popularity bias.

In some popular music streaming platforms, such as Spotify, users might
experience a lack of discoverability, where they might be recommended a set of
songs that are popular or belong to the same artists, which might narrow the
user's music discovery experience and harm the smaller artists and independent
artists.In \cite{dugeri2024legal}, the author, using Spotify as a studying
case, mentions that the top 43,000 artists, roughly 1.4\% of the total aritsts
on the platform, pull in 90\% of the royalties.

To the extent of my knowledge, there has not been much work being done using
topic modeling for music recommendation. However, I want to attribute the
inspiration for this work to the paper \cite{tsukuda2017lyricjumper}, which
uses a similar approach to recommend songs based on the lyrics, whose work
apply the LDA model and modify the model to create their own model, which is
called the ArtistTaste model(AT), which belives that the artists might have
many topics (themes) and each song will be associated with one topic.

\section{Topic Modeling}
\subsection{Introduction}
Topic modeling is a statistical learning technique that aims to discover
abstract "topics" or thematic structures within a collection of documents.
Unlike supervised classification, topic modeling operates in an unsupervised
manner, automatically identifying latent semantic patterns without requiring
labeled data or predefined categories.

At its core, topic modeling assumes that documents are generated from a mixture
of topics, where each topic represents a probability distribution over words in
the vocabulary. This generative assumption allows the model to decompose
complex textual data into interpretable thematic components. For instance, in a
corpus of news articles, topics might emerge as "politics," "sports," or
"technology," each characterized by words like \{election, vote, government\},
\{game, team, score\}, or \{software, computer, algorithm\}, respectively.

The mathematical foundation of topic modeling rests on the concept of
dimensionality reduction. Given a corpus of $D$ documents with vocabulary size
$V$, traditional bag-of-words representations result in high-dimensional,
sparse feature spaces. Topic modeling projects this high-dimensional space into
a lower-dimensional topic space of size $K$ (where $K \ll V$), capturing the
most salient semantic patterns while reducing noise and computational
complexity.

In the context of music recommendation, topic modeling offers several
advantages over traditional approaches. First, it provides interpretable
insights into the thematic content of songs, allowing us to understand why
certain recommendations are made. Second, it naturally handles the
heterogeneity of musical content, as songs can contain multiple themes (e.g.,
love and loss, celebration and nostalgia). Third, as mentoioned before,
utilizing only the lyrics aspect, though simple as it might sound, reduces
dependence on external metadata or collaborative signals, making it
particularly valuable for cold-start scenarios or underrepresented artists.

The evolution of topic modeling has progressed from classical probabilistic
methods to neural approaches that leverage modern deep learning techniques.
Classical models like Latent Dirichlet Allocation (LDA) \cite{blei2003latent}
employ Bayesian inference to estimate topic distributions, while neural topic
models such as Contextualized Topic Models (CTM) \cite{bianchi2021cross}
integrate pre-trained embeddings to capture semantic relationships more
effectively. This progression has led to improved topic coherence, better
handling of short texts, and enhanced cross-lingual capabilities. Another model
examined in this work is BERTopic \cite{grootendorst2022bertopic}, which
combines transformer embeddings, dimensionality reduction algorithm as well as
with clustering algorithms to generate interpretable topics.

\subsection{Problem Formulation}
To formalize the topic modeling problem, we begin by defining the fundamental
components of textual data representation. In topic modeling, we work with
three hierarchical levels of data organization:

\textbf{Words} constitute the atomic units of our analysis. Each word is drawn from a fixed vocabulary $\mathcal{V} = \{1, 2, \ldots, V\}$ of size $V$. Following standard practice in natural language processing, we represent words using one-hot encoding vectors. Specifically, the $v$-th word in the vocabulary is represented as a $V$-dimensional binary vector $\mathbf{w}$ where the $v$-th component equals 1 and all other components are 0. This sparse representation allows us to uniquely identify each word while maintaining computational efficiency.

\textbf{Documents} are sequences of words that form coherent textual units. In our context, each document corresponds to a song's lyrics. Formally, a document is represented as $\mathbf{w} = (w_1, w_2, \ldots, w_N)$, where $N$ denotes the document length and $w_n$ represents the $n$-th word in the sequence. Unlike bag-of-words models that ignore word order, this sequential representation preserves the structural information within lyrics, though topic models typically aggregate these sequences into frequency counts.

\textbf{Corpora} represent collections of documents that share thematic or contextual relationships. Our corpus $\mathcal{D} = \{\mathbf{w}_1, \mathbf{w}_2, \ldots, \mathbf{w}_M\}$ consists of $M$ documents (song lyrics), where each document $\mathbf{w}_m$ belongs to our music dataset. The corpus serves as the foundation for discovering latent topical patterns that span across individual songs and artists.

The central challenge in topic modeling lies in discovering latent thematic
structures within this hierarchical representation. Given only the observed
words across documents, we aim to infer hidden topics that explain the
statistical patterns in word co-occurrence. This unsupervised learning problem
requires models that can capture both the diversity of themes within documents
and the coherence of themes across the corpus.

\subsection{Latent Dirichlet Allocation (LDA)}
Latent Dirichlet Allocation (LDA) \cite{blei2003latent} is a widely used
probabilistic topic model that represents documents as mixtures of topics. The
fundamental assumption of LDA is that each document is generated through a
two-stage generative process. First, a document-specific mixture of topics is
drawn from a Dirichlet distribution with parameter $\boldsymbol{\alpha}$.
Second, for each word position in the document, a topic is selected according
to this mixture, and then a word is drawn from the chosen topic's word
distribution.

Formally, LDA assumes the following generative model for a corpus of $M$
documents:

\begin{enumerate}
    \item For each topic $k \in \{1, \ldots, K\}$:
          \begin{itemize}
              \item Draw topic-word distribution $\boldsymbol{\beta}_k \sim
                        \text{Dir}(\boldsymbol{\eta})$
          \end{itemize}
    \item For each document $m \in \{1, \ldots, M\}$:
          \begin{itemize}
              \item Draw document-topic distribution $\boldsymbol{\theta}_m \sim
                        \text{Dir}(\boldsymbol{\alpha})$
              \item For each word position $n \in \{1, \ldots, N_m\}$:
                    \begin{itemize}
                        \item Draw topic assignment $z_{m,n} \sim \text{Multinomial}(\boldsymbol{\theta}_m)$
                        \item Draw word $w_{m,n} \sim \text{Multinomial}(\boldsymbol{\beta}_{z_{m,n}})$
                    \end{itemize}
          \end{itemize}
\end{enumerate}

The key parameters are:
\begin{itemize}
    \item $\boldsymbol{\alpha}$: Dirichlet prior on document-topic distributions (controls topic sparsity per document)
    \item $\boldsymbol{\eta}$: Dirichlet prior on topic-word distributions (controls word sparsity per topic)
    \item $K$: Number of topics (hyperparameter)
\end{itemize}

The joint probability of observed words and latent variables is:
\begin{equation}
    p(\mathbf{w}, \mathbf{z}, \boldsymbol{\theta}, \boldsymbol{\beta} | \boldsymbol{\alpha}, \boldsymbol{\eta}) = \prod_{k=1}^K p(\boldsymbol{\beta}_k | \boldsymbol{\eta}) \prod_{m=1}^M p(\boldsymbol{\theta}_m | \boldsymbol{\alpha}) \prod_{n=1}^{N_m} p(z_{m,n} | \boldsymbol{\theta}_m) p(w_{m,n} | \boldsymbol{\beta}_{z_{m,n}})
\end{equation}

Inference in LDA typically employs variational Bayes or Gibbs sampling to
estimate the posterior distributions of topics and document-topic mixtures. The
model's strength lies in its interpretability: topics emerge as coherent word
clusters, and documents are represented as weighted combinations of these
thematic components.

\begin{figure}[htbp]
    \centering
    \includegraphics[width=0.6\textwidth]{Figures/lda_graphical_model.png}
    \caption{Graphical model representation of Latent Dirichlet Allocation (LDA). The plate notation shows the generative process where $\alpha$ and $\beta$ are hyperparameters, $\theta$ represents document-topic distributions, $z$ denotes topic assignments for words, and $w$ represents observed words. The inner plate indicates $N$ words per document, while the outer plate represents $M$ documents in the corpus.}
    \label{fig:lda_model}
\end{figure}

\subsection{BERTopic}
BERTopic \cite{grootendorst2022bertopic} represents a modern approach to topic
modeling that leverages advances in transformer-based language models and
clustering algorithms. Unlike classical probabilistic methods such as LDA,
BERTopic adopts a three-stage pipeline that combines dense embeddings,
dimensionality reduction, and clustering to discover topics in textual data.

\subsubsection{Model Architecture}
The BERTopic framework consists of three main components working in sequence:

\textbf{Document Embedding Generation:} BERTopic begins by converting documents into dense vector representations using pre-trained transformer models. By default, it employs sentence-transformers models (such as \texttt{all-MiniLM-L6-v2}) that have been fine-tuned to produce semantically meaningful embeddings. For a document $d$, this stage produces an embedding vector $\mathbf{e}_d \in \mathbb{R}^{H}$, where $H$ is the embedding dimension (typically 384 or 768). These embeddings capture semantic relationships between documents more effectively than traditional bag-of-words representations.

\textbf{Dimensionality Reduction:} The high-dimensional embeddings are then reduced to a lower-dimensional space using UMAP (Uniform Manifold Approximation and Projection). This step serves two purposes: computational efficiency and improved clustering performance. UMAP preserves both local and global structure in the data, making it particularly suitable for subsequent clustering. The dimensionality reduction can be formalized as a mapping $f: \mathbb{R}^{H} \rightarrow \mathbb{R}^{d}$ where $d \ll H$.

\textbf{Clustering:} The reduced embeddings are clustered using HDBSCAN (Hierarchical Density-Based Spatial Clustering of Applications with Noise), which automatically determines the number of clusters and can identify outliers. Unlike k-means clustering, HDBSCAN does not require specifying the number of clusters a priori and can discover clusters of varying densities and shapes. Documents that cannot be assigned to any cluster are labeled as noise/outliers.

\subsubsection{Topic Representation}
After clustering, BERTopic generates interpretable topic representations using
a class-based TF-IDF approach. For each cluster $c$, the method:

\begin{enumerate}
    \item Concatenates all documents belonging to cluster $c$ into a single
          pseudo-document
    \item Calculates TF-IDF scores, treating each cluster as a separate class
    \item Extracts the top-k words with highest TF-IDF scores as the topic representation
\end{enumerate}

The class-based TF-IDF modification enhances the discriminative power of words
within topics by considering the frequency of words across different clusters.
This results in topic representations that better capture the unique
characteristics of each cluster.

\subsubsection{Advantages and Limitations}
BERTopic offers several advantages over traditional topic models:
\begin{itemize}
    \item \textbf{Semantic Understanding:} Transformer embeddings capture contextual word meanings, leading to more coherent topics
    \item \textbf{Automatic Topic Detection:} No need to specify the number of topics beforehand
    \item \textbf{Outlier Handling:} HDBSCAN naturally identifies and separates noise from meaningful clusters
    \item \textbf{Scalability:} Efficient processing of large document collections
\end{itemize}

However, BERTopic also has notable limitations:
\begin{itemize}
    \item \textbf{Hard Assignment:} Each document belongs to exactly one topic, limiting the model's ability to capture multi-thematic content
    \item \textbf{Deterministic Clustering:} Unlike probabilistic models, BERTopic does not provide uncertainty estimates or topic distributions
    \item \textbf{Limited Interpretability:} The embedding and clustering stages are less interpretable than the probabilistic generative process in LDA
    \item \textbf{Hyperparameter Sensitivity:} Performance can be sensitive to UMAP and HDBSCAN parameters
\end{itemize}

\subsection{Product of Experts Latent Dirichlet Allocation (ProdLDA)}
Product of Experts Latent Dirichlet Allocation (ProdLDA)
\cite{srivastava2017autoencoding} represents a crucial bridge between classical
probabilistic topic models and modern neural approaches. Developed as part of
the Autoencoding Variational Inference framework, ProdLDA addresses several
limitations of traditional LDA while maintaining its probabilistic foundation,
making it an essential stepping stone toward more advanced neural topic models
like CTM.

\subsubsection{Motivation and Background}
Traditional LDA inference methods, such as Gibbs sampling and variational
inference, face several computational and modeling challenges:

\begin{itemize}
    \item \textbf{Scalability Issues:} Standard inference algorithms can be computationally expensive for large corpora
    \item \textbf{Local Optima:} Variational inference often gets trapped in poor local optima
    \item \textbf{Limited Expressiveness:} The Dirichlet prior constrains the model's flexibility in capturing complex topic structures
    \item \textbf{Difficulty with Sparse Data:} Traditional methods struggle with short documents or sparse word co-occurrence patterns
\end{itemize}

ProdLDA addresses these issues by reformulating topic modeling within a neural
variational autoencoder framework, enabling more flexible inference while
preserving the interpretable probabilistic structure of topic models.

\subsubsection{Model Architecture}
ProdLDA replaces the traditional Dirichlet prior over document-topic
distributions with a logistic-normal distribution, which can be more easily
handled by neural networks. The key architectural components include:

\textbf{Encoder Network:} The encoder maps the bag-of-words representation of a document $\mathbf{w}_d$ to the parameters of a multivariate Gaussian distribution in an unconstrained space. Specifically, it learns:
\begin{equation}
    q(\mathbf{h}_d | \mathbf{w}_d) = \mathcal{N}(\boldsymbol{\mu}_d, \text{diag}(\boldsymbol{\sigma}_d^2))
\end{equation}
where $\boldsymbol{\mu}_d$ and $\boldsymbol{\sigma}_d$ are outputs of neural networks taking $\mathbf{w}_d$ as input.

\textbf{Logistic-Normal Transform:} The latent representation $\mathbf{h}_d$ is transformed to the probability simplex using the softmax function:
\begin{equation}
    \boldsymbol{\theta}_d = \text{softmax}(\mathbf{h}_d)
\end{equation}
This ensures that $\boldsymbol{\theta}_d$ represents a valid probability distribution over topics.

\textbf{Decoder Network:} The decoder reconstructs the document using the topic distribution. Unlike traditional LDA where topic-word distributions are drawn from Dirichlet priors, ProdLDA parameterizes these distributions using neural networks:
\begin{equation}
    p(\mathbf{w}_d | \boldsymbol{\theta}_d) = \text{Multinomial}(\text{softmax}(\boldsymbol{\theta}_d^T \mathbf{B}))
\end{equation}
where $\mathbf{B} \in \mathbb{R}^{K \times V}$ is a learned topic-word matrix.

\subsubsection{Training Objective}
ProdLDA optimizes the Evidence Lower BOund (ELBO) of the log-likelihood:
\begin{equation}
    \mathcal{L} = \mathbb{E}_{q(\mathbf{h}_d|\mathbf{w}_d)}[\log p(\mathbf{w}_d|\mathbf{h}_d)] - \text{KL}(q(\mathbf{h}_d|\mathbf{w}_d) \| p(\mathbf{h}_d))
\end{equation}

The first term encourages accurate reconstruction of documents, while the KL
divergence term acts as a regularizer, preventing the latent representations
from deviating too far from the prior distribution $p(\mathbf{h}_d) =
    \mathcal{N}(0, I)$.

\subsubsection{Advantages over Traditional LDA}
ProdLDA offers several improvements over classical LDA:

\begin{itemize}
    \item \textbf{Scalable Inference:} The neural architecture enables efficient gradient-based optimization and mini-batch training
    \item \textbf{Better Topic Quality:} The logistic-normal distribution often produces more coherent topics compared to Dirichlet priors
    \item \textbf{Flexible Architecture:} Neural networks allow for more complex relationships between documents and topics
    \item \textbf{Handling Sparse Data:} The continuous latent space better handles documents with few words
\end{itemize}

\subsubsection{Limitations and Considerations}
Despite its advantages, ProdLDA also has limitations:

\begin{itemize}
    \item \textbf{Loss of Interpretability:} The neural parameterization makes the model less interpretable than classical LDA
    \item \textbf{Hyperparameter Sensitivity:} Performance depends on architectural choices and learning rate tuning
    \item \textbf{Computational Requirements:} Requires more computational resources than traditional inference methods
    \item \textbf{Limited Contextual Understanding:} Still relies on bag-of-words representations, missing semantic relationships
\end{itemize}

\subsubsection{Foundation for CTM}
ProdLDA serves as the crucial foundation for Contextualized Topic Models (CTM).
While ProdLDA maintains the bag-of-words input limitation, it establishes the
neural variational framework that CTM extends by incorporating pre-trained
contextual embeddings. The architectural innovations in ProdLDA—particularly
the use of neural networks for parameterizing topic distributions and the
logistic-normal latent space—directly enable CTM's ability to leverage
transformer-based representations while maintaining probabilistic topic
modeling principles.

\subsection{Contextualized Topic Models (CTM)}
Contextualized Topic Models (CTM) \cite{bianchi2021cross} represent a direct
and elegant extension of ProdLDA that addresses its primary limitation: the
reliance on bag-of-words representations. While ProdLDA established an
effective neural variational framework for topic modeling, it still suffered
from the semantic limitations inherent in traditional bag-of-words document
representations. CTM solves this by replacing bag-of-words inputs with rich,
contextualized document embeddings while maintaining the exact same
probabilistic framework and training objectives as ProdLDA.

\subsubsection{Core Innovation: Contextualized Document Representations}
The fundamental innovation of CTM lies in its input representation strategy.
Instead of using bag-of-words vectors $\mathbf{w}_d$ as in ProdLDA, CTM employs
contextualized document embeddings $\mathbf{x}_d$ generated by pre-trained
transformer models.

\textbf{Document Encoding Process:} For a document $d$ containing words $\{w_1, w_2, \ldots, w_N\}$, CTM:
\begin{enumerate}
    \item Computes contextualized word embeddings using pre-trained transformers (BERT,
          RoBERTa, or sentence-transformers)
    \item Aggregates these embeddings (typically through averaging or pooling) to form a
          document-level representation $\mathbf{x}_d \in \mathbb{R}^H$
    \item Uses $\mathbf{x}_d$ as input to the ProdLDA architecture instead of the
          bag-of-words vector $\mathbf{w}_d$
\end{enumerate}

This substitution enables CTM to capture semantic relationships that
bag-of-words representations cannot express. For instance, documents containing
synonymous phrases like "happy song" and "joyful melody" will have similar
embeddings $\mathbf{x}_d$, leading to similar topic distributions, whereas
their bag-of-words representations would be orthogonal.

\subsubsection{Architectural Framework}
CTM adopts the identical neural variational autoencoder architecture as
ProdLDA, with only the input representation modified:

\textbf{Encoder Network:} The encoder maps contextualized document representations $\mathbf{x}_d$ (instead of bag-of-words $\mathbf{w}_d$) to Gaussian parameters:
\begin{equation}
    q(\mathbf{h}_d | \mathbf{x}_d) = \mathcal{N}(\boldsymbol{\mu}_d, \text{diag}(\boldsymbol{\sigma}_d^2))
\end{equation}
where $\boldsymbol{\mu}_d$ and $\boldsymbol{\sigma}_d$ are neural network outputs taking the contextualized embedding $\mathbf{x}_d$ as input.

\textbf{Logistic-Normal Transform and Decoder:} The remaining architecture is identical to ProdLDA:
\begin{align}
    \boldsymbol{\theta}_d                   & = \text{softmax}(\mathbf{h}_d)                                           \\
    p(\mathbf{w}_d | \boldsymbol{\theta}_d) & = \text{Multinomial}(\text{softmax}(\boldsymbol{\theta}_d^T \mathbf{B}))
\end{align}

\textbf{Training Objective:} CTM optimizes the same ELBO as ProdLDA, but with contextualized inputs:
\begin{equation}
    \mathcal{L} = \mathbb{E}_{q(\mathbf{h}_d|\mathbf{x}_d)}[\log p(\mathbf{w}_d|\mathbf{h}_d)] - \text{KL}(q(\mathbf{h}_d|\mathbf{x}_d) \| p(\mathbf{h}_d))
\end{equation}

The key insight is that while the decoder still reconstructs bag-of-words
representations (to maintain topic interpretability), the encoder leverages
rich semantic information from contextualized embeddings.

\subsubsection{Advantages Through Contextualization}
The contextualized input representation provides CTM with several advantages
over ProdLDA:

\begin{itemize}
    \item \textbf{Semantic Understanding:} Documents with similar meanings but different vocabularies are mapped to similar topic distributions
    \item \textbf{Better Topic Coherence:} Topics emerge based on semantic similarity rather than just word co-occurrence patterns
    \item \textbf{Improved Short Text Handling:} Contextualized embeddings provide richer representations for documents with few words
    \item \textbf{Cross-lingual Capabilities:} When using multilingual embeddings, topics can span across languages
    \item \textbf{Reduced Vocabulary Dependence:} The model is less sensitive to specific word choices and can generalize better to unseen vocabulary
\end{itemize}

\subsubsection{Implementation Considerations}
The practical implementation of CTM requires careful consideration of the
embedding generation process:

\textbf{Embedding Model Selection:} The choice of pre-trained transformer affects topic quality. Sentence-transformers models like \texttt{all-MiniLM-L6-v2} are commonly used for their balance of semantic quality and computational efficiency.

\textbf{Computational Trade-offs:} While embedding generation adds computational overhead, the improved semantic representations often justify this cost, especially for applications requiring high topic coherence.

\subsubsection{Relationship to ProdLDA}
CTM can be viewed as ProdLDA with a sophisticated preprocessing step that
replaces bag-of-words vectors with contextualized embeddings. This perspective
highlights that:

\begin{itemize}
    \item The core probabilistic framework remains unchanged
    \item All mathematical properties of ProdLDA (convergence guarantees, interpretation
          of latent variables) apply to CTM
    \item The improvement in performance comes primarily from better input
          representations rather than architectural innovations
    \item CTM inherits ProdLDA's advantages (scalable inference, flexible neural
          architecture) while addressing its semantic limitations
\end{itemize}

\subsubsection{Application to Music Lyrics}
For music recommendation based on lyrics, CTM's contextualized approach is
particularly valuable:

\textbf{Capturing Lyrical Nuance:} Song lyrics often use metaphorical language, wordplay, and cultural references that contextualized embeddings can better represent than bag-of-words approaches.

\textbf{Handling Artistic Diversity:} Different artists may express similar themes using vastly different vocabulary. CTM can identify these semantic similarities and group songs with related themes regardless of specific word choices.

\textbf{Genre-Independent Themes:} Contextualized representations enable CTM to discover thematic patterns that transcend genre boundaries, potentially revealing interesting cross-genre recommendations based on lyrical content.

\section{Methodology}

Now, I will present the methodology for this work, which can be divided into 4
steps: dataset construction, lyrics extraction and preprocessing, model
training, and finally recommendation generation. We will now walk through each
step in detail.

\subsection{Dataset Construction and Lyrics Extraction}
By using ChatGPT to generate a list of the most representative artists across
various genres that span through the decades in the US-UK music, I assemble a
dataset of 600 artists. Most of these artists have songs written in the English
language; however, there are some artists that have songs written in other
languages, such as Spanish, Korean, and Japanese, though the majority of the
songs are in English. Then, using the LyricsGenius Python library
\cite{miller2019lyricsgenius}, which provides a convenient interface to the
Genius.com API, we extract the top 10 songs for each artist, resulting in
approximately 6,000 song lyrics documents. After extraction,

\subsection{Preprocessing}
After the lyrics extraction, I then move on to perform some preprocessing steps
before feeding the lyrics into the topic modeling algorithm. For example, a song's lyrics might look like this:
\begin{verbatim}
    I love you, baby
    And if it's quite alright
    I need you, baby
    To warm the lonely night
\end{verbatim}




\subsection{Model Training}
We trained LDA (\texttt{gensim}), BERTopic, and CTM (via \texttt{OCTIS}) with
$K=10$ topics. Hyperparameters were tuned on a held-out validation set.

\subsection{Artist-level Aggregation}
Document-topic distributions ($\theta_d$) were averaged per artist to form
artist-topic profiles ($\Theta_a = \frac{1}{|D_a|} \sum_{d\in D_a} \theta_d$),
reflecting an artist's thematic signature.

\subsection{Recommendation Generation}
Given a query artist or song, we computed cosine similarity between profiles.
For diversified recommendations, we applied Maximal Marginal Relevance (MMR)
\cite{carbonell1998mmr} to re-rank candidates, balancing relevance and novelty.

\section{Evaluation}
\subsection{Topic Coherence and Diversity}
We measured coherence (UMass) and topic diversity \cite{dieng2020evaluating}.

\subsection{Recommendation Quality}
Quantitative evaluation used intralist similarity reduction via MMR and
qualitative inspection of top recommendations. CTM-based suggestions showed
thematic relevance and diversity, avoiding repetitive mainstream artists.

\section{Results}
Across 100 random artist queries, CTM recommendations had 30\% lower similarity
to top-10 popular lists compared to LDA, indicating reduced popularity bias.
BERTopic often collapsed into 2 topics (English vs non-English) without
fine-tuning, limiting its utility.

\section{Discussion}
Our findings demonstrate that contextualized topic models effectively capture
lyrical themes, producing coherent and diverse artist recommendations.
Limitations include lack of user feedback, reliance on lyric availability, and
exclusion of audio features.

\section{Conclusion and Future Work}
We presented a lyrics-based music recommendation pipeline leveraging topic
modeling. Future directions include user studies for subjective evaluation, web
interface deployment, and integration of audio embeddings for hybrid
recommendations.

% \bibliographystyle{plainnat}
\bibliographystyle{unsrt}
\bibliography{references}

\end{document}
